\section{Evaluation}
\label{sec:eval}

The evaluation is a measurement of the model in Section~\ref{sec:model} and of the tracker in Section~\ref{sec:design}.
No end-to-end number from the two preceding studies is carried forward.
Those studies used different budgets (a 48\,GB residency target and a 20\,GB cap), different model sets, and different primary metrics.
Filling a table with either set would answer a different question than P1--P4.
This section therefore specifies the protocol, the hypotheses, and the tables.
Cells that require the new runs are marked unfilled.

\subsection{Hypotheses}
\label{sec:hyp}

Each hypothesis names the prediction it tests and the comparison that rejects it.
\begin{description}
  \item[H1 (P1).] On every model--device pair with \(\PhiRatio \ge 1\), the exposed expert-transfer latency at \(S=0\) is no higher than at any \(S>0\) under the same pin quota. Rejection: some positive \(S\) reduces p95 exposed wait by a margin larger than the run-to-run standard deviation.
  \item[H2 (P2).] On every pair with \(\PhiRatio < 1\) and a non-empty feasible set, the measured latency-minimizing horizon is non-increasing in \(C\) and non-decreasing in prefetch size \(m\). Rejection: the minimizing \(S\) is unchanged, within one layer, when the prefetched representation switches between the two tiers or when the same binary is moved between the A6000 and H20.
  \item[H3 (P3).] At a fixed budget with \(\PhiRatio < 1\), the feasible set is empty on the prefill of the batch-32 workload in Table~\ref{tab:density} and non-empty on the following decode step, for at least one model. Rejection: the feasible-set predicate has the same truth value on those two steps for every model in the sweep.
  \item[H4 (P4).] The largest \(n_{\mathrm{hi}}\) that preserves a non-empty feasible set equals the slack bound from~\eqref{eq:slack}, to within one expert per layer. In the dense sub-phase that bound is zero. Rejection: adding a pin beyond the bound leaves p95 wait unchanged, or withholding a pin inside the bound reduces wait.
  \item[H5 (proxy).] With the latency band held, hotness-ranked pins improve task score over a uniform low-precision assignment that uses the same \(B\). Rejection: the task-score confidence intervals overlap on every model, or the pin policy meets the score gain only by leaving the latency band.
  \item[H6 (coupling).] The naive composition, which computes a pin quota and a horizon independently against the full budget \(B\), either exceeds \(B\) or has higher p95 exposed wait than \systemname\ at the same budget. Rejection: the composition stays inside \(B\) and matches \systemname\ within run-to-run noise on every phase.
\end{description}
H2 is the existence condition for a joint paper.
If the minimizing horizon does not move when \(m\) changes, precision and prefetch are sequential optimizations and a single budget controller is unnecessary.
The sweep reports that result even if it rejects H2.

\subsection{Systems Under Test}
\label{sec:systems}

All systems share the dispatcher, the pools, the NVML sampler, and the quantization recipe.
They differ only in the policy that chooses the published tier and the horizon.
\begin{description}
  \item[Uniform.] Every expert is resident at the highest uniform tier whose pool fits in \(B\). No online changes.
  \item[Static mixed.] The Phase~II quota is computed from \(B\), identities are the top of a calibration ranking, and both the quota and the identities are frozen. This is the offline mixed-precision assignment~\cite{duanmu2025mxmoe,chitty2025mopeq}.
  \item[Precision-only.] Phase~II rules at every budget: full low-precision floor when it fits, hotness pins, \(S=0\). When the floor does not fit, the policy is infeasible and the run is recorded as infeasible rather than quietly quantized further.
  \item[Horizon-only.] Phase~III rules with \(n_{\mathrm{hi}}=0\) and a single tier, the largest tier whose working set the horizon controller keeps inside \(B\). This is the adaptive-horizon policy with fixed expert size~\cite{song2024promoe,shen2025expertflow}.
  \item[Miss-path mixed.] Horizon-only, except a miss is fetched at low precision and promoted only after a second hit. This isolates the ``cheapen the miss'' policy~\cite{tang2024hobbit} from a pin that is protected across periods.
  \item[Naive composition.] Precision-only and horizon-only run together. Each computes its set against \(B\). Admission is still on, so the second reservation fails instead of overflowing; we count failed reservations and the resulting misses.
  \item[\systemname.] Algorithm~\ref{alg:period}.
  \item[Oracle horizon.] For each workload, an offline sweep of \(S\) and \(n_{\mathrm{hi}}\) selects the pair with the lowest p95 exposed wait subject to bytes \(\le B\). The oracle is a ceiling for the latency model, not a deployable policy. It is computed on the calibration prompts and also, separately, on the test prompts. The gap between those two oracles is the price of non-stationarity.
\end{description}
We do not report a comparison against a system we could not execute on the same checkpoint.
Where a published runtime does not load one of the checkpoints, that cell is absent, not approximated.

\subsection{Testbed, Models, and Workloads}
\label{sec:testbed}

Devices are an NVIDIA RTX~A6000 (48\,GB), an NVIDIA~H20, and an Ascend~910B, the three machines on which the preliminary horizon sweep was taken.
The A6000 is the only device on which the full quality suite is required.
H20 and 910B are required for H2, because they change \(C\).
If the 910B software stack cannot execute the INT4 kernels for a checkpoint, that pair is dropped and the drop is reported; it is not replaced with a different quantization.

Models are Qwen3-30B-A3B, Qwen3-Next-80B-A3B, and Phi-3.5-MoE, plus DeepSeek-V2-Lite, Qwen1.5-MoE, and Qwen2-MoE so that the 20\,GB region of the sweep contains models whose routing we have already traced~\cite{qwen3technicalreport,qwen3nextmodelcard,abdin2024phi3,bi2024deepseek,yang2024qwen2}.
Table~\ref{tab:models} records the structural parameters.
The budget sweep is \(B \in \{12, 16, 20, 24, 32, 40, 48\}\)\,GB on the A6000, clipped to the device capacity on the others, and further clipped to \(P_{\mathrm{hi}}\) because budgets above Phase~I are idle.
For each pair we report \(\PhiRatio\) from the measured \(P_{\mathrm{lo}}\), not from the nominal parameter count.

\begin{table}[t]
  \centering
  \caption{Model structure. Activated parameters are the vendor figures. \(\PhiRatio\) is filled from measured \(P_{\mathrm{lo}}\) at each \(B\), not from this table.}
  \label{tab:models}
  \small
  \begin{tabular}{@{}lrrrrr@{}}
    \toprule
    Model & Params & Active & Layers & Experts & Top-\(k\) \\
    \midrule
    Qwen3-30B-A3B & 30.5B & 3.3B & 48 & 128 & 8 \\
    Qwen3-Next-80B-A3B & 80B & 3B & 48 & 512 & 10 \\
    Phi-3.5-MoE & 42B & 6.6B & 32 & 16 & 2 \\
    DeepSeek-V2-Lite & \tbd{fill} & \tbd{fill} & \tbd{fill} & \tbd{fill} & \tbd{fill} \\
    Qwen1.5-MoE & \tbd{fill} & \tbd{fill} & \tbd{fill} & \tbd{fill} & \tbd{fill} \\
    Qwen2-MoE & \tbd{fill} & \tbd{fill} & \tbd{fill} & \tbd{fill} & \tbd{fill} \\
    \bottomrule
  \end{tabular}
\end{table}

Latency requests are drawn from ShareGPT~\cite{ShareGPT_Vicuna_unfiltered}, bucketed by prompt length, with at most 50 requests per bucket and length variation held within 5\% inside a bucket.
This is a public conversational corpus; the runs send no data to an external service.
Each latency configuration uses five warmup iterations and 100 measured iterations, repeated for five seeds.
We report the mean, the median, and the 95th and 99th percentiles of exposed wait, of miss-path time, of time to first token, and of time per output token.
Percentiles are computed on the pooled raw samples of a seed and then summarized across seeds by the median and the range.
Exposed wait is the time a layer is stalled on a missing expert after subtracting the overlap with that layer’s resident expert compute.
Miss-path time is the migration-stream duration and is allowed to be larger than exposed wait.

Quality uses WikiText-2 perplexity~\cite{merity2016pointer} and the task scores of MMLU-Pro, GPQA, GSM8K, and HumanEval~\cite{wang2024mmlu,rein2024gpqa,cobbe2021gsm8k,chen2021evaluating}.
The item sets are the pinned subsets already used in our routing traces: 200 MMLU-Pro items, the GPQA-Diamond set, 200 GSM8K items, and the full HumanEval set.
Perplexity uses at most 128 non-overlapping windows of 2{,}048 tokens.
Multiple-choice items are scored by the conditional log-likelihood of the full answer label.
HumanEval is pass@1 under greedy decoding after executing the official tests.
Accuracy intervals are Wilson intervals at 95\%.
We do not pool tasks into a single significance test.
A model is a separate experiment.

Calibration and test are disjoint.
Hotness initialization, the frozen controller constants, and the oracle-horizon selection use a calibration slice of ShareGPT and a calibration slice of WikiText.
Test prompts, including the three task workloads used for the hot-set disjointness check, are held out.
Constants are not refit per device.
\(C\) and \(T_{\ell}\) are measured online and are not calibrated constants.

\subsection{Experiments}
\label{sec:experiments}

\paragraph{E0: repeat the phenomena on one stack.}
Before any policy comparison we remeasure Table~\ref{tab:density}, the Layer-15 top-10 overlap across WikiText, GSM8K, and HumanEval, and the latency-versus-\(S\) curve at \(B=20\)\,GB.
E0 uses the uniform low tier and a fixed horizon, with \systemname’s loops off.
If the unique-expert ratio does not move by at least 30 percentage points from prefill to decode on Qwen3-30B at batch 32, or if the latency curve is not U-shaped on at least two models, the preliminary motivation does not survive the unified runtime and the later experiments are not informative.
E0 is reported in full, including a negative outcome.

\paragraph{E1: the phase diagram.}
For each model on the A6000 we plot p95 exposed wait and WikiText perplexity against \(B\), with one curve per system in Section~\ref{sec:systems}.
The horizontal axis is annotated with the measured \(\PhiRatio=1\) and \(\PhiRatio=P_{\mathrm{hi}}/P_{\mathrm{lo}}\) points.
H1 requires the precision-only curve and the \systemname\ curve to meet for \(\PhiRatio\ge 1\), and the horizon-only curve to lie above them by the migration it still performs.
H6 is read from the same figure: the naive composition either has no point, because admission refused the dual reservation, or lies above \systemname.
Table~\ref{tab:phase} is the numerical companion, filled per regime rather than by a single speedup.

\begin{table}[t]
  \centering
  \caption{Phase diagram on the A6000, p95 exposed wait (ms) and perplexity.
  One row group per model.
  Unfilled cells are the joint re-measurement.}
  \label{tab:phase}
  \small
  \begin{tabular}{@{}llrrrr@{}}
    \toprule
    Model & Regime & Uniform & Prec.-only & Horiz.-only & \systemname \\
    \midrule
    \tbd{model} & \(\PhiRatio \ge \rho\) & \tbd{ms / ppl} & \tbd{ms / ppl} & \tbd{ms / ppl} & \tbd{ms / ppl} \\
    \tbd{model} & \(1 \le \PhiRatio < \rho\) & \tbd{ms / ppl} & \tbd{ms / ppl} & \tbd{ms / ppl} & \tbd{ms / ppl} \\
    \tbd{model} & \(\PhiRatio < 1\), sparse & \tbd{ms / ppl} & \tbd{infeasible or ms} & \tbd{ms / ppl} & \tbd{ms / ppl} \\
    \tbd{model} & \(\PhiRatio < 1\), dense & \tbd{ms / ppl} & \tbd{infeasible or ms} & \tbd{ms / ppl} & \tbd{ms / ppl} \\
    \bottomrule
  \end{tabular}
\end{table}

\paragraph{E2: does \(S^{\star}\) move with \(m\) and with \(C\)?}
On each Phase~III configuration we sweep \(S\) with prefetch size fixed at \(m^{\mathrm{lo}}\) and again at \(m^{\mathrm{hi}}\), pins disabled, and we record the minimizing \(S\).
We repeat the low-precision sweep on H20 and on the 910B.
H2 is a statement about the argmin, not about the height of the curve.
The model error is \(|S^{\star}_{\mathrm{pred}}-S^{\star}_{\mathrm{meas}}|\), with \(S^{\star}_{\mathrm{pred}}\) computed from~\eqref{eq:sstar} twice: once with idle \(C\), once with in-overlap \(C\).
We report the mean absolute error in layers.
A useful model is one whose in-overlap error is at most one layer on a majority of configurations.
An idle-\(C\) error that is large while the in-overlap error is small is evidence for the contention caveat in Section~\ref{sec:model-limits}, not evidence against the overlap identity.

\paragraph{E3: inside one request.}
For the batch-32 ShareGPT buckets we log, per request, the feasible-set predicate on the last prefill layer and on each of the first 32 decode steps, together with the applied \(S\) and \(n_{\mathrm{hi}}\).
H3 is a predicate on that log.
We also report the number of periods the applied \(S\) takes to come within one layer of a new \(S^{\star}\) after the prefill-to-decode boundary, and the migration bytes moved during those periods.
The hysteresis and the \(\pm 1\) step are acceptable only if that migration stays below the migration of an ablation that sets \(S\) to \(S^{\star}\) in one period.

\paragraph{E4: slack and pins.}
At three budgets in Phase~III (sparse) we compare \(n_{\mathrm{hi}} \in \{0, n^{\star}, n^{\star}{+}1\}\), where \(n^{\star}\) is the slack bound.
H4 predicts that \(n^{\star}{+}1\) either fails admission or increases p95 wait, and that \(n^{\star}\) does not.
The dense rows of the same table use the batch-32 prefill.
H4 predicts \(n^{\star}=0\) there.

\paragraph{E5: quality at a fixed latency band.}
On the A6000 we take the largest \(B\) in Phase~II and the largest \(B\) in Phase~III at which \systemname’s p95 time-to-first-token stays within 10\% of the uniform-low baseline’s.
If no Phase~III budget meets that band, the Phase~III quality comparison is reported as ``band infeasible'' and H5 is not claimed there.
Inside the band we report task scores for uniform, static mixed, precision-only, and \systemname.
Table~\ref{tab:quality} has one block per model.
The static-mixed column is the comparison that tests whether online identity changes matter once the quota is held fixed.
Intervals that overlap are reported as overlaps.

\begin{table}[t]
  \centering
  \caption{Task score at the latency band of Section~\ref{sec:experiments}, A6000.
  Wilson 95\% intervals travel with each accuracy cell in the artifact.}
  \label{tab:quality}
  \small
  \begin{tabular}{@{}llrrrr@{}}
    \toprule
    Model & Task & Uniform & Static mixed & Prec.-only & \systemname \\
    \midrule
    \tbd{model} & MMLU-Pro & \tbd{score} & \tbd{score} & \tbd{score} & \tbd{score} \\
    \tbd{model} & GPQA & \tbd{score} & \tbd{score} & \tbd{score} & \tbd{score} \\
    \tbd{model} & GSM8K & \tbd{score} & \tbd{score} & \tbd{score} & \tbd{score} \\
    \tbd{model} & HumanEval & \tbd{score} & \tbd{score} & \tbd{score} & \tbd{score} \\
    \tbd{model} & WikiText ppl & \tbd{ppl} & \tbd{ppl} & \tbd{ppl} & \tbd{ppl} \\
    \bottomrule
  \end{tabular}
\end{table}

\paragraph{E6: invariant, overhead, and ablations.}
For every reported run we record the maximum of the runtime byte counter and the NVML process peak.
A counter above \(B\) is a failed implementation.
A process peak above the device, while the counter stays inside \(B\), is a reserve violation and is reported separately.
Controller CPU time per fast period and per slow period is reported as a distribution.
The design target is that the fast period stays well below \(T_{\ell}\); the evaluation states the measured fraction of \(T_{\ell}\) rather than a microsecond budget chosen in advance.

Ablations remove one mechanism at a time on the A6000, on one Phase~II model and one Phase~III model: no hysteresis on pins, no \(\pm 1\) limit, idle \(C\) instead of in-overlap \(C\), mean miss size instead of p95, no cache-aware issue order, and a single pool instead of per-tier pools.
A further ablation replaces the pregate-or-frequency predictor with a CPU random forest trained on calibration tuples of token identity, layer, and horizon.
That ablation is informative only when E2 shows a residual wait that the overlap model attributes to miss-set error rather than to \(C\) or \(m\).
If the residual is small, the forest is not part of the system we recommend.

\subsection{What We Will Not Claim from This Protocol}
\label{sec:nonclaims}

A single speedup against uniform quantization is not a result under this protocol.
Uniform quantization is the Phase~I optimum and a feasible Phase~II baseline; it is infeasible in part of Phase~III, and a speedup there would compare against a system that does not fit.
A task-score gain with an unstated latency regression is not evidence for H5.
A horizon policy that wins at one batch size and one device is the behavior P2 predicts for a fixed \(S\), and it is not a competing system result unless the policy is the horizon-only controller, which already adapts \(S\).
